\chapter{What monetary signals can and cannot tell us}\label{ch:money}
\question{Why can a financially attractive action have harmful consequences outside its price?}

The clinic's invoice is not an illusion. Payment lets a supplier meet obligations, purchase replacement parts and pay workers. Prices allow many unlike goods to be compared within a common monetary account. A household can make a budget; a firm can examine revenue and expense; a public provider can plan expenditure. These are practical achievements.

A price is also a signal formed within institutions. It can change when production becomes harder, demand rises, alternatives disappear, contracts change or market power shifts. Its meaning depends on who can offer, who can buy and what rules apply. A signal produced by these conditions need not be a complete description of the physical consequences of an action.

\section{Purchasing power and credit}
Purchasing power concerns what a person or institution can obtain with available monetary resources at prevailing prices. It is different from need. Two people may need the same essential service while having very different capacities to pay. A price records the terms of an exchange; it does not directly measure the urgency of every excluded person's need.

Credit connects present spending to future obligations. In its explanation of modern UK banking, the Bank of England describes how a commercial-bank loan creates a matching deposit. It also explains constraints on lending, including profitability, regulation and monetary policy. This is an institutional mechanism, not unlimited creation of physical wealth.\cite{boe}

Imagine a loan financing the clinic's pump. It can make an order possible today, ahead of future revenues. The pump still requires manufacture, transport, installation and maintenance. Creating a claim and supplying the physical capacity are connected activities, but only the second installs equipment. Conversely, refusing a loan can leave usable technical possibilities unrealised. The physical and financial questions remain linked without becoming identical.

\section{Consequences outside the transaction}
An externality is an effect on people outside a transaction that is not fully reflected in the decision-makers' private costs or benefits. Public goods pose a related problem when people cannot readily be excluded from the benefit and one person's use does not subtract from another's. These concepts explain why private incentives can differ from wider consequences; they do not say that every market always fails.\cite{externalities}

Consider a hypothetical workshop choosing whether to contain a harmful discharge. If its invoice includes the cost of disposal but excludes damage downstream, the cheapest option for the workshop can burden others. The example assumes an unpriced effect. It does not establish how common that situation is or measure its total cost.

Change the rules and the decision can change. A discharge limit, an enforceable liability rule or a charge tied to emissions can make prevention part of the workshop's own calculation. Public provision can organise a service that individual purchases would undersupply. Rules need information and enforcement, but their existence matters to the argument.

The EU Emissions Trading System supplies an actual example of a connection between a physical quantity and a monetary incentive: covered operators monitor emissions and surrender allowances within a cap-and-trade system. This does not establish that the instrument covers every ecological effect or that its price is adequate for every purpose. It does refute the claim that environmental effects never enter monetary institutions.\cite{ets}

\section{A price and a physical history}
Suppose two deliveries cost the clinic the same amount. One uses a well-maintained pipe, the other an emergency tanker. Equal payments do not logically imply equal fuel use, water loss, reliability or exposure to disruption. To answer those questions, we need measurements and a boundary for the comparison.

The reverse also holds. Two physically similar deliveries may have different prices because ownership, taxes, contracts or bargaining conditions differ. A physical account does not determine those legal terms by itself. Neither price nor physical description has to be declared meaningless for their difference to matter.

There are deliberate feedback institutions already. The ECB's 2025 strategy names price stability as its primary objective and commits to considering climate and nature implications within its mandate. A stated monetary mandate is not a guarantee of every social or ecological function, but it is plainly a feedback arrangement.\cite{ecb}

For EBU, the question is narrower than whether money is good or bad. Could a declared physical-action account make certain consequences explicit in a way that is useful, auditable and sufficiently local? It would have to earn its place alongside existing accounts and rules. Calling a number physical does not automatically make it complete or well calibrated.

\practice{A pollution charge makes the workshop's cleaner option cheaper. Does this make a physical record redundant?}{No. The charge itself depends on a defined effect and some means of measurement or estimation. The record and incentive can support each other. Nor does the charge's existence prove all downstream effects have been included.}

\carry{Omitted effects and inability to purchase essentials are different problems. The next chapter asks why a resource can exist while a person's access to it remains inadequate.}
